\documentclass[10pt,a4paper]{amsart}
\usepackage{amsmath,amssymb,amsthm,graphicx}
% Retained preamble for the Stacks Project chapter snapshot in this directory.
% The chapter's own preamble.tex is a copy of the shared Stacks Project preamble; the block
% below is sliced MECHANICALLY, line for line, from it. Nothing was added, removed or
% redefined; the wrapper class, the xr-hyper/hyperref wiring and the externaldocument list
% are omitted because the assembler supplies the document class.
%
% lines taken from preamble.tex: 3, 5, 9, 11, 12, 156, 157, 158, ... (51 lines total)
%
% package coverage restored from the chapter's own preamble (2026-10-13): xy, xy-2cell

\usepackage[all]{xy}
\xyoption{2cell}
\UseAllTwocells

\newenvironment{reference}{\comment}{\endcomment}
\newenvironment{slogan}{\comment}{\endcomment}
\newenvironment{history}{\comment}{\endcomment}
\theoremstyle{plain}
\newtheorem{theorem}[subsection]{Theorem}
\newtheorem{proposition}[subsection]{Proposition}
\newtheorem{lemma}[subsection]{Lemma}
\theoremstyle{definition}
\newtheorem{definition}[subsection]{Definition}
\newtheorem{example}[subsection]{Example}
\newtheorem{exercise}[subsection]{Exercise}
\newtheorem{situation}[subsection]{Situation}
\theoremstyle{remark}
\newtheorem{remark}[subsection]{Remark}
\newtheorem{remarks}[subsection]{Remarks}
\numberwithin{equation}{subsection}
\def\lim{\mathop{\mathrm{lim}}\nolimits}
\def\colim{\mathop{\mathrm{colim}}\nolimits}
\def\Spec{\mathop{\mathrm{Spec}}}
\def\Hom{\mathop{\mathrm{Hom}}\nolimits}
\def\Ext{\mathop{\mathrm{Ext}}\nolimits}
\def\SheafHom{\mathop{\mathcal{H}\!\mathit{om}}\nolimits}
\def\SheafExt{\mathop{\mathcal{E}\!\mathit{xt}}\nolimits}
\def\Sch{\mathit{Sch}}
\def\Mor{\mathop{\mathrm{Mor}}\nolimits}
\def\Ob{\mathop{\mathrm{Ob}}\nolimits}
\def\Sh{\mathop{\mathit{Sh}}\nolimits}
\def\NL{\mathop{N\!L}\nolimits}
\def\CH{\mathop{\mathrm{CH}}\nolimits}
\def\proetale{{pro\text{-}\acute{e}tale}}
\def\etale{{\acute{e}tale}}
\def\QCoh{\mathit{QCoh}}
\def\Ker{\mathop{\mathrm{Ker}}}
\def\Im{\mathop{\mathrm{Im}}}
\def\Coker{\mathop{\mathrm{Coker}}}
\def\Coim{\mathop{\mathrm{Coim}}}
\DeclareMathSymbol{\boxtimes}{\mathbin}{AMSa}{"02}
\def\QCohstack{\mathcal{QC}\!\mathit{oh}}
\def\Cohstack{\mathcal{C}\!\mathit{oh}}
\def\Spacesstack{\mathcal{S}\!\mathit{paces}}
\def\Quotfunctor{\mathrm{Quot}}
\def\Hilbfunctor{\mathrm{Hilb}}
\def\Curvesstack{\mathcal{C}\!\mathit{urves}}
\def\Polarizedstack{\mathcal{P}\!\mathit{olarized}}
\def\Complexesstack{\mathcal{C}\!\mathit{omplexes}}
\def\Pic{\mathop{\mathrm{Pic}}\nolimits}
\def\Picardstack{\mathcal{P}\!\mathit{ic}}
\def\Picardfunctor{\mathrm{Pic}}
\def\Deformationcategory{\mathcal{D}\!\mathit{ef}}
\usepackage{cleveref}
\crefname{originallemma}{original lemma}{original lemmas}
\Crefname{originallemma}{Original lemma}{Original lemmas}
\makeatletter
\newlabel{lemma-curves-qs-lfp}{{0DSS}{source}{}{sourcepointer-1}{}}
\newlabel{lemma-curves-qs-lfp@cref}{{[originallemma][0][]0DSS}{[1][0][]source}{}{}{}}
\makeatother
\crefname{originallemma}{original lemma}{original lemmas}
\Crefname{originallemma}{Original lemma}{Original lemmas}
\makeatletter
\newlabel{more-morphisms-lemma-deform-property}{{06AG}{source}{}{sourcepointer-2}{}}
\newlabel{more-morphisms-lemma-deform-property@cref}{{[originallemma][0][]06AG}{[1][0][]source}{}{}{}}
\makeatother
\crefname{originallemma}{original lemma}{original lemmas}
\Crefname{originallemma}{Original lemma}{Original lemmas}
\makeatletter
\newlabel{spaces-over-fields-lemma-codim-1-point-in-schematic-locus}{{0ADD}{source}{}{sourcepointer-3}{}}
\newlabel{spaces-over-fields-lemma-codim-1-point-in-schematic-locus@cref}{{[originallemma][0][]0ADD}{[1][0][]source}{}{}{}}
\makeatother
\crefname{originallemma}{original lemma}{original lemmas}
\Crefname{originallemma}{Original lemma}{Original lemmas}
\makeatletter
\newlabel{stacks-geometry-lemma-characterize-smoothness}{{0DZS}{source}{}{sourcepointer-4}{}}
\newlabel{stacks-geometry-lemma-characterize-smoothness@cref}{{[originallemma][0][]0DZS}{[1][0][]source}{}{}{}}
\makeatother
\crefname{originallemma}{original lemma}{original lemmas}
\Crefname{originallemma}{Original lemma}{Original lemmas}
\makeatletter
\newlabel{stacks-morphisms-lemma-where-smooth}{{0DZR}{source}{}{sourcepointer-5}{}}
\newlabel{stacks-morphisms-lemma-where-smooth@cref}{{[originallemma][0][]0DZR}{[1][0][]source}{}{}{}}
\makeatother
\title{Stacks, Moduli of Curves: lemma-in-smooth-locus}
\author{The Stacks Project Authors}
\date{stacks-project:moduli-curves.tex:local-snapshot-20261012; GNU FDL 1.2 or later}
\begin{document}
\maketitle
\noindent\textit{Editorial note.} Statement span is lines 1626-1634 of sources/stacks-moduli-curves/pinned-original-source.tex; proof span is lines 1636-1678, the complete author proof (41 body lines).\par\medskip
\noindent\textit{Editorial note.} Navigation stubs: 5 cited-but-not-included labels are supplied as non-load-bearing pointers with their pinned tags where the snapshot tag index has them.\par\medskip
\noindent\textit{Editorial note.} extra\_packages is empty on purpose: the retained preamble already loads xy with options and a second declaration clashes in the assembler wrapper.\par\medskip
\noindent\textit{Editorial note.} Text is retained verbatim, including the author's own typos and placeholders.\par\medskip
\noindent\textit{Editorial note.} Queue index 8 of queue\_moduli-curves.json (pb=2035, stmt 1626-1634, proof 1636-1678).\par\medskip
\noindent\textit{Allowed standard prerequisites (editorial).}\begin{itemize}
\item Explicit prerequisite boundary. The cited inputs are prerequisites; the argument in the retained span is not replaced by any cited result.
\end{itemize}\medskip
\ifdefined\hypertarget\hypertarget{sourcepointer-1}{}\fi
\noindent\textit{Source reference.} Original lemma 0DSS: sources/stacks-moduli-curves/pinned-original-source.tex, source lines 655--658. Cited by the retained statement or proof but not part of the collected spans; supplied only as a navigation pointer. The argument does not depend on its content.\par\medskip
\ifdefined\hypertarget\hypertarget{sourcepointer-2}{}\fi
\noindent\textit{Source reference.} Original lemma 06AG: sources/stacks-moduli-curves/pinned-original-source.tex, source lines 1678--1678. Cited by the retained statement or proof but not part of the collected spans; supplied only as a navigation pointer. The argument does not depend on its content.\par\medskip
\ifdefined\hypertarget\hypertarget{sourcepointer-3}{}\fi
\noindent\textit{Source reference.} Original lemma 0ADD: sources/stacks-moduli-curves/pinned-original-source.tex, source lines 1678--1678. Cited by the retained statement or proof but not part of the collected spans; supplied only as a navigation pointer. The argument does not depend on its content.\par\medskip
\ifdefined\hypertarget\hypertarget{sourcepointer-4}{}\fi
\noindent\textit{Source reference.} Original lemma 0DZS: sources/stacks-moduli-curves/pinned-original-source.tex, source lines 1678--1678. Cited by the retained statement or proof but not part of the collected spans; supplied only as a navigation pointer. The argument does not depend on its content.\par\medskip
\ifdefined\hypertarget\hypertarget{sourcepointer-5}{}\fi
\noindent\textit{Source reference.} Original lemma 0DZR: sources/stacks-moduli-curves/pinned-original-source.tex, source lines 1678--1678. Cited by the retained statement or proof but not part of the collected spans; supplied only as a navigation pointer. The argument does not depend on its content.\par\medskip
\begin{lemma}
\label{lemma-in-smooth-locus}
In the situation above the following are equivalent
\begin{enumerate}
\item the classifying morphism $\Spec(k) \to \Curvesstack$ factors
through the open where $\Curvesstack \to \Spec(\mathbf{Z})$ is smooth,
\item the deformation category $\Deformationcategory_X$ is unobstructed.
\end{enumerate}
\end{lemma}

\begin{proof}
Since $\Curvesstack \longrightarrow \Spec(\mathbf{Z})$ is locally
of finite presentation (Lemma \ref{lemma-curves-qs-lfp})
formation of the open substack where
$\Curvesstack \longrightarrow \Spec(\mathbf{Z})$ is smooth commutes with
flat base change
(Morphisms of Stacks, Lemma \ref{stacks-morphisms-lemma-where-smooth}).
Since the Cohen ring $\Lambda$ is flat over $\mathbf{Z}$,
we may work over $\Lambda$. In other words, we are trying to prove that
$$
\Lambda\text{-}\Curvesstack \longrightarrow \Spec(\Lambda)
$$
is smooth in an open neighbourhood of the point
$x_0 : \Spec(k) \to \Lambda\text{-}\Curvesstack$
defined by $X/k$ if and only if $\Deformationcategory_X$ is unobstructed.

\medskip\noindent
The lemma now follows from
Geometry of Stacks, Lemma \ref{stacks-geometry-lemma-characterize-smoothness}
and the equality
$$
\Deformationcategory_X =
\mathcal{F}_{\Lambda\text{-}\Curvesstack, k, x_0}
$$
This equality is not completely trivial to establish. Namely, on the left
hand side we have the deformation category classifying all flat deformations
$Y \to \Spec(A)$ of $X$ as a scheme over $A \in \Ob(\mathcal{C}_\Lambda)$.
On the right hand side we have the deformation category classifying all
flat morphisms $Y \to \Spec(A)$ with special fibre $X$
where $Y$ is an algebraic space and
$Y \to \Spec(A)$ is proper, of finite presentation, and of
relative dimension $\leq 1$. Since $A$ is Artinian, we find
that $Y$ is a scheme for example by Spaces over Fields, Lemma
\ref{spaces-over-fields-lemma-codim-1-point-in-schematic-locus}.
Thus it remains to show: a flat deformation $Y \to \Spec(A)$ of
$X$ as a scheme over an Artinian local ring $A$ with residue field $k$
is proper, of finite presentation, and of relative dimension $\leq 1$.
Relative dimension is defined in terms of fibres and hence holds
automatically for $Y/A$ since it holds for $X/k$.
The morphism $Y \to \Spec(A)$ is proper and locally of finite presentation
as this is true for $X \to \Spec(k)$, see
More on Morphisms, Lemma \ref{more-morphisms-lemma-deform-property}.
\end{proof}

\end{document}
